\chapter{El primer arranque}\label{cap:05-primer-arranque}

\section{Qué vas a ver}

Al encender aparece el logo de Lázaro con el pulso y, después, la pantalla
para entrar: elige tu usuario y escribe la contraseña. Si marcaste
\emph{Iniciar sesión automáticamente}, entra solo al escritorio.

\begin{nota}
Si en el disco también está Windows, al encender aparece primero un menú
negro, llamado \termino{GRUB}. La primera línea arranca Lázaro (se elige sola
a los pocos segundos); para Windows, baja con las flechas hasta
\emph{Windows Boot Manager} y pulsa \keys{Enter}.
\end{nota}

\section{Lo que pasa solo en los primeros minutos}

\begin{itemize}
  \item \textbf{La hora.} Si aparece corrida, se corrige sola en unos minutos
        al conectarse a internet.
  \item \textbf{Las aplicaciones de Flathub} que marcaste (OnlyOffice,
        Chrome, Brave) se descargan en segundo plano cuando hay internet.
        Dales unos minutos; aparecerán en el menú.
  \item \textbf{Las herramientas de Lázaro} se conectan a su repositorio para
        recibir mejoras con las actualizaciones normales
        (capítulo~\ref{cap:15-herramientas}).
\end{itemize}

\section{Si instalaste la Ligera sin internet}\label{sec:primer-arranque-texto}

Con la edición Ligera, una parte de la configuración final necesita
internet. Si instalaste sin conexión, el equipo arranca en una
\textbf{pantalla de texto} con un aviso. No es un error:

\begin{pasos}
  \item Conecta el cable de red, si tienes. Si no, entra con tu usuario y
        contraseña y conecta el wifi escribiendo:
\begin{terminal}
nmcli device wifi connect "NOMBRE DE LA RED" --ask
\end{terminal}
        Te pide la clave del wifi. Mientras la escribes no se ve nada en la
        pantalla: es normal.
  \item Espera. En unos minutos la configuración termina sola y aparece la
        pantalla de inicio. No hace falta reiniciar.
\end{pasos}

\section{Primeros pasos recomendados}

\begin{pasos}
  \item \textbf{Conecta el wifi} desde el icono de red de la barra
        (capítulo~\ref{cap:07-internet}). Se recuerda para siempre.
  \item \textbf{Actualiza.} Si aparece el icono de actualizaciones en la
        barra, haz clic y sigue los pasos (sección~\ref{sec:actualizar}).
  \item \textbf{Prueba la impresora}, si tienes (capítulo~\ref{cap:12-imprimir}).
  \item \textbf{Trae tus archivos} de la copia que hiciste
        (capítulo~\ref{cap:11-archivos}).
  \item \textbf{Abre Firefox} e inicia sesión en tu correo y tus cuentas.
  \item \textbf{Hojea la Parte II} de este manual para conocer lo que
        trae Lázaro.
\end{pasos}

El \textbf{Manual de Lázaro} queda siempre en el menú de aplicaciones, por
si lo necesitas.

\fuentes
